% Technical validation of the Phase 2 models
% Numbers come from reports/dl/checks/*.csv (python reports/dl/checks/dl_checks.py validate);
% the runs and checks were made on one NVIDIA T4 GPU (Kaggle).

\section{Technical Validation}
\label{sec:dl_validation}

A single test score can mislead in five ways: another seed may give a different model, the test fires may
by chance favour one model, the probabilities may not mean what they say, the model may need more data, and
test patches may repeat training locations. We check each of these for the final U-Net (focal loss), the CNN
(weighted BCE) and the Phase~1 XGBoost, refitted with its committed hyperparameters to obtain its test
predictions (PR AUC 0.2576 against the committed 0.2565, Appendix~\ref{app:checks}). All PR AUC values are on
the official test split, with pixels labelled $-1$ excluded.

\subsection{Seeds}

Each model was trained with seeds 0, 1 and 2 (for XGBoost the seed sets row and column subsampling).
Table~\ref{tab:dl_seeds} gives the mean and standard deviation. The U-Net is ahead on every metric. The two
neural networks vary four to five times more between seeds than XGBoost (standard deviation 0.008 to 0.010
in PR AUC against 0.002), so a single run of a network can be off by about 0.01.

\begin{table*}[t]
\centering
\caption{Test results over three seeds (mean $\pm$ standard deviation).}
\label{tab:dl_seeds}
\small
\setlength{\tabcolsep}{6pt}
\renewcommand{\arraystretch}{1.12}
\begin{tabular}{lccccc}
\toprule
\rowcolor{tblhead}
\textbf{Model} & \textbf{PR AUC} & \textbf{ROC AUC} & \textbf{F1} & \textbf{R@1\%} & \textbf{R@5\%} \\
\midrule
\rowcolor{tblours}
U-Net (focal) & \textbf{0.275} $\pm$ .008 & \textbf{0.923} $\pm$ .003 & \textbf{0.349} $\pm$ .006 & \textbf{0.398} $\pm$ .004 & \textbf{0.711} $\pm$ .005 \\
CNN (BCE) & 0.249 $\pm$ .010 & 0.878 $\pm$ .011 & 0.335 $\pm$ .008 & 0.378 $\pm$ .009 & 0.653 $\pm$ .006 \\
XGBoost & 0.255 $\pm$ .002 & 0.905 $\pm$ .001 & 0.330 $\pm$ .000 & 0.388 $\pm$ .001 & 0.673 $\pm$ .001 \\
\bottomrule
\end{tabular}
\end{table*}

\subsection{Reproducibility}

We retrained the seed-0 U-Nets and CNN a second time with the same seed, settings and GPU. No rerun
reproduced its predictions: the mean change per pixel was 0.03 to 0.07, the kept epoch moved by up to four
epochs, and test PR AUC changed by $-0.022$ to $+0.013$. GPU kernels that add numbers in a varying order make
training non-deterministic, so on a GPU the same seed does not give the same model, and here the spread of
reruns is as large as the spread of seeds. Such run-to-run variance can be as large as the differences between
methods being compared~\cite{bouthillier2021}. We therefore report means over seeds and paired
comparisons rather than single runs. Across hardware, the same configuration agrees within this noise: the
BCE U-Net reaches 0.269 $\pm$ 0.010 on the T4 against 0.269 for the reference run on an Apple GPU. Bit-exact
reruns would need deterministic algorithms in PyTorch, at a cost in speed.

\subsection{Bootstrap confidence intervals}

As in Phase~1 (Appendix~\ref{app:checks}), we resample the test set 1\,000 times~\cite{efron1979}, either by
patch or by whole fire event (the 257 events of Section~\ref{sec:events}), and recompute PR AUC each time.
Table~\ref{tab:dl_ci} and Fig.~\ref{fig:dl_ci} give the intervals for the seed-0 models. As in Phase~1, the
event intervals are wide (about $\pm$0.06), because a few large fires hold most of the fire pixels, and
those of the three models overlap.

\begin{table}[t]
\centering
\caption{Test PR AUC of the seed-0 models with 95\% bootstrap intervals (1\,000 resamples).}
\label{tab:dl_ci}
\small
\renewcommand{\arraystretch}{1.12}
\begin{tabular}{lccc}
\toprule
\rowcolor{tblhead}
\textbf{Model} & \textbf{PR AUC} & \textbf{Patches} & \textbf{Events} \\
\midrule
\rowcolor{tblours}
U-Net (focal) & 0.2694 & 0.253-0.286 & 0.211-0.329 \\
CNN (BCE) & 0.2591 & 0.241-0.277 & 0.180-0.333 \\
XGBoost & 0.2576 & 0.242-0.273 & 0.190-0.323 \\
Persistence & 0.1065 & 0.097-0.117 & 0.059-0.161 \\
\bottomrule
\end{tabular}
\end{table}

\begin{figure}[t]
\centering
\includegraphics[width=0.85\columnwidth]{dl/figures/dl_ci_events.png}
\caption{Test PR AUC with 95\% intervals over fire events (seed 0); diamonds are the three seeds.}
\label{fig:dl_ci}
\end{figure}

Paired intervals are much narrower, because both models are scored on the same resampled fires. Compared
with XGBoost (seed 0) on the same event resamples, every U-Net seed is better: by 0.012 (0.001 to 0.025)
for seed 0, 0.014 (0.002 to 0.028) for seed 1 and 0.027 (0.015 to 0.039) for seed 2. The CNN is not: its
differences are $+0.001$ ($-0.011$ to 0.013), $-0.009$ ($-0.022$ to 0.005) and $-0.018$ ($-0.030$ to
$-0.008$). The gain comes with the U-Net's architecture rather than with the loss or with neural networks as
such: trained with the same BCE loss, the U-Net still reaches 0.269 $\pm$ 0.010 against 0.249 $\pm$ 0.010
for the CNN. The U-Net sees a $44\times44$-pixel window around each pixel, the CNN only $9\times9$, which is
about the local context of the Phase~1 neighbourhood features, and the CNN scores about the same as XGBoost.

\subsection{Calibration}

Training weights fire pixels by $w = 90.0$ and keeps every pixel, so by the prior-shift argument of
Section~\ref{sec:theory_imb} (with $r = 1$) the raw scores overstate the odds of fire by about 90, and we
undo this with (\ref{eq:inverse_shift}). Table~\ref{tab:dl_cal} gives the Brier score~\cite{brier1950}, the mean
squared difference between probability and outcome, and the Brier skill score against always predicting
the training fire rate (1.10\%).

\begin{table}[t]
\centering
\caption{Calibration on test (seed 0): Brier score of the raw and corrected probabilities, skill score
(BSS) and expected calibration error (ECE) over deciles of the corrected probability.}
\label{tab:dl_cal}
\small
\setlength{\tabcolsep}{3.5pt}
\renewcommand{\arraystretch}{1.12}
\begin{tabular}{lccccc}
\toprule
\rowcolor{tblhead}
\textbf{Model} & \textbf{Raw} & \textbf{Corrected} & \textbf{BSS} & \textbf{ECE} & \textbf{Mean $p$} \\
\midrule
U-Net (focal) & 0.1067 & 0.0119 & 0.043 & 0.0090 & 0.0061 \\
CNN (BCE) & 0.0807 & 0.0104 & 0.159 & 0.0029 & 0.0113 \\
XGBoost & 0.0856 & 0.0104 & 0.159 & 0.0019 & 0.0107 \\
\midrule
\multicolumn{6}{l}{Always 1.10\%: Brier 0.0124; observed fire rate on test 0.0125.} \\
\bottomrule
\end{tabular}
\end{table}

\begin{figure}[t]
\centering
\includegraphics[width=0.62\columnwidth]{dl/figures/dl_calibration.png}
\caption{Reliability of the corrected probabilities on test: mean probability against observed fire rate
in each tenth of the pixels (log scales).}
\label{fig:dl_cal}
\end{figure}

The raw scores are badly calibrated for every model, with a Brier score six to nine times worse than the
constant forecast. The correction fixes this for the two models trained with weighted BCE: the CNN and
XGBoost reach a skill score of 0.16, and their top tenth of pixels predicts 9.5\% and 8.9\% fire against
8.7\% and 9.7\% observed (Fig.~\ref{fig:dl_cal}). Focal loss breaks the correction, because its scores are
less extreme (Section~\ref{sec:dl_loss}; see also~\cite{mukhoti2020}). After the correction, the U-Net's
probabilities are too flat: 0.15\% in the lowest tenth against 0.06\% observed, and 2.5\% in the top tenth
against 10.0\% observed, for a skill score of only 0.04. Its ranking, and so its PR AUC, is unaffected. If
the U-Net's probabilities are to be read as fire risk, they need a recalibration fitted on eval, such as
Platt scaling.

\subsection{Learning curve}

We trained the focal U-Net (seed 0) on 5, 10, 25 and 50\% of the training patches. Each run had the same
budget of gradient steps and the same patience in steps as the full run (at most $30/f$ epochs, patience
$5/f$, for a share $f$). Unlike XGBoost, whose eval PR AUC stays between 0.204 and 0.207 from 5 to 100\% of
the data (Section~\ref{sec:theory_bv}), the U-Net uses more data (Fig.~\ref{fig:dl_lc}): its eval PR AUC
rises from 0.194 at 5\% to 0.224 at 50\% and is flat afterwards (0.223 at 100\%, within the rerun noise).
With 10\% of the training patches it already matches XGBoost trained on all of them (0.211 against 0.207).
The U-Net learns spatial features that the per-pixel model could not, but beyond half of the training set
more data of the same kind adds little.

\begin{figure}[t]
\centering
\includegraphics[width=0.75\columnwidth]{dl/figures/dl_learning_curve.png}
\caption{Eval PR AUC against the share of training patches: focal U-Net (one run per share) and the Phase~1
XGBoost.}
\label{fig:dl_lc}
\end{figure}

\subsection{Location-overlap split}

503 of the 1\,689 test patches (30\%) cover exactly the same location as a training patch
(Section~\ref{sec:method}). Scored separately (Table~\ref{tab:dl_overlap}), these seen locations are easier
for every model, persistence included. On the 1\,186 new locations alone, the U-Net still beats XGBoost by
0.013 (0.004 to 0.026, paired over fire events), while the CNN and XGBoost remain level ($+0.001$, from
$-0.014$ to 0.013). The U-Net's lead is therefore not a product of locations it has seen before.

\begin{table}[t]
\centering
\caption{Test PR AUC on all, new and seen locations (seed 0).}
\label{tab:dl_overlap}
\small
\renewcommand{\arraystretch}{1.12}
\begin{tabular}{lccc}
\toprule
\rowcolor{tblhead}
\textbf{Model} & \textbf{All (1\,689)} & \textbf{New (1\,186)} & \textbf{Seen (503)} \\
\midrule
\rowcolor{tblours}
U-Net (focal) & 0.269 & 0.255 & 0.335 \\
CNN (BCE) & 0.259 & 0.243 & 0.330 \\
XGBoost & 0.258 & 0.242 & 0.324 \\
Persistence & 0.107 & 0.097 & 0.151 \\
\bottomrule
\end{tabular}
\end{table}

\begin{projectbox}
What the checks support: the U-Net with focal loss is the best model. Its lead over XGBoost, 0.01 to 0.03 in
PR AUC, holds for each of three seeds when both are scored on the same fires, and for seed 0 also on
unseen locations. The CNN is level with XGBoost. Network training on a GPU is not reproducible run by run, so
results are means over seeds. After the closed-form correction, the probabilities of the weighted-BCE models
are close to real risk; those of the focal U-Net need a recalibration before they are used as risk.
\end{projectbox}
